%! suppress = Makeatletter
%! suppress = TooLargeSection
%! suppress = MissingLabel
\documentclass{article}

% Fields
\usepackage{geometry}
\geometry{top=25mm}
\geometry{bottom=35mm}
\geometry{left=20mm}
\geometry{right=20mm}
% ------------------------------------------------

% Graphics
\usepackage{color}
\usepackage{tabularx}
\usepackage{tikz}
\usepackage{blkarray}
\usepackage{graphicx}
% ------------------------------------------------

% Math
\usepackage{amsmath, amsfonts}
\usepackage{amssymb}
\usepackage{proof}
\usepackage{mathrsfs}
% Crossed-out symbols
% https://tex.stackexchange.com/questions/75525/how-to-write-crossed-out-math-in-latex
\usepackage[makeroom]{cancel}
\usepackage{mathtools}
% ------------------------------------------------

% Additional font sizes
% https://www.overleaf.com/learn/latex/Questions/How_do_I_adjust_the_font_size%3F
\usepackage{moresize}
% Additional colors
% https://www.overleaf.com/learn/latex/Using_colours_in_LaTeX
\usepackage{xcolor}
% \texttimes
\usepackage{textcomp}
% ------------------------------------------------

% Language
\usepackage[utf8] {inputenc}
\usepackage[T2A] {fontenc}
\usepackage[english, russian] {babel}
\usepackage{indentfirst, verbatim}
\usetikzlibrary{cd, babel}
% ------------------------------------------------

% Fonts
\usepackage{stmaryrd}
\usepackage{cmbright}
\usepackage{wasysym}
% ------------------------------------------------

% Code
% https://tex.stackexchange.com/questions/99475/how-to-invoke-latex-with-the-shell-escape-flag-in-texstudio-former-texmakerx
% Colored, requires --shell-escape compiling option
\usepackage{minted}
\setminted{xleftmargin=\parindent, autogobble, escapeinside=\#\#}
\usepackage{listings}
% ------------------------------------------------

% Custom envs
% https://tex.stackexchange.com/questions/371286/draw-a-horizontal-line-in-latex
\newenvironment{proof}{\subparagraph{\hspace{-1em}Решение:\newline}\-}{\par\noindent\rule{\textwidth}{0.4pt}}
% ------------------------------------------------

% Custom commands
\newcommand{\comb}[1]{\mathbf{#1}}
\newcommand{\step}{\rightsquigarrow}
\newcommand{\term}[1]{\mathbf{#1}}
\newcommand{\ap}{~}
\newcommand{\termdef}{\coloneqq}
\newcommand{\subst}[3]{\left[#2 \mapsto #3 \right] #1}
\newcommand{\eqbeta}{=_\beta}
\newcommand{\eqbetaeta}{=_{\beta\eta}}
\newcommand{\betastep}{\rightarrow_\beta}
\newcommand{\betasteps}{\twoheadrightarrow_\beta}
\newcommand{\eqeta}{=_\eta}
\newcommand{\tlist}[1]{\term{[}#1\term{]}} % list-term
\renewcommand{\emph}[1]{{\color{blue} \textit{#1}}}
% ------------------------------------------------

% Head
\usepackage{fancyhdr}
\usepackage{hyperref}
\pagestyle{fancy}
\fancyhead[R]{Ярошевский Илья}
\fancyhead[L]{ИТМО MSE, ФП 2024, Дз 3}
% ------------------------------------------------

% Numbering
% https://tex.stackexchange.com/questions/80113/hide-section-numbers-but-keep-numbering
\makeatletter
\renewcommand\thesubsection{Блок \@arabic\c@subsection.\hspace{-0.8em}}
\renewcommand\thesubsubsection{Задание \@arabic\c@subsection.\@arabic\c@subsubsection\hspace{-0.8em}}
% https://tex.stackexchange.com/questions/327689/numbering-subsubsections-with-letters
\renewcommand\theparagraph{\alph{paragraph})\hspace{-0.8em}}
% https://tex.stackexchange.com/questions/129208/numbering-paragraphs-in-latex
\setcounter{secnumdepth}{4}
\makeatother
% ------------------------------------------------

\begin{document}
    \section*{Дз 3. Просто-типизированное лямбда-исчисление}

    Глобальный минимум: 8б.

    \subsection{Нахождение наиболее общего типа}

    Локальный минимум: 2б.

    Найдите наиболее общий тип.
    Например, у $\term{K}$ наиболее общий тип $\alpha \to \beta \to \alpha$: мы знаем, что тип первого аргумента такой же, как и тип возвращаемого значения.
    $\term{K} : \alpha \to \alpha \to \alpha$~--- это корректная типизация
    для терма $\term{K}$, но не наиболее общая (она накладывает дополнительное необоснованное ограничение на совпадение типов аргументов).
    Её можно получить из предыдущей подстановкой $\subst{}{\beta}{\alpha}$.

    \subsubsection{(1б)}

    \paragraph{} $\lambda x\ldotp x$

    \begin{proof}
        \begin{gather}
            \lambda x^{?_0} \ldotp x \\
            \tau = ?_0 \to ?_0
        \end{gather}
        $\lambda x \ldotp x : \alpha \to \alpha$
    \end{proof}

    \paragraph{} $\lambda x~x\ldotp x$

    \begin{proof}
        \begin{gather}
            \lambda x^{?_0}~x^{?_1}\ldotp x \\
            \tau = ?_0 \to ?_1 \to ?_1
        \end{gather}
        $\lambda x~x\ldotp x : \alpha \to \beta \to \beta$
    \end{proof}

    \paragraph{} $\lambda x^{\alpha \to \alpha}\ldotp x$

    \begin{proof}
        \begin{gather}
            \lambda x^{\alpha \to \alpha}\ldotp x \\
            \tau = (\alpha \to \alpha) \to (\alpha \to \alpha)
        \end{gather}
        $\lambda x^{\alpha \to \alpha}\ldotp x : (\alpha \to \alpha) \to (\alpha \to \alpha)$
    \end{proof}

    \paragraph{} $\lambda x~y\ldotp x \ap y$

    \begin{proof}
        \begin{gather}
            \lambda x^{?_0}~y^{?_1}\ldotp \underbrace{x~y}_{?_2} \\
            \begin{cases}
                ?_0 = ?_1 \to ?_2 \\
                \tau = ?_0 \to ?_1 \to ?_2
            \end{cases}
            \begin{cases}
                \tau = (?_1 \to ?_2) \to ?_1 \to ?_2
            \end{cases}
        \end{gather}
        $\lambda x~y\ldotp x~y : (\alpha \to \beta) \to \alpha \to \beta$
    \end{proof}

    \subsubsection{(1б)}

    \paragraph{} $\lambda x\ldotp x \ap \term{K} \ap \term{I} \ap x$

    \begin{proof}
        \begin{gather}
            \lambda x^{?_0}\ldotp \underbrace{\underbrace{\underbrace{x^{?_0} \ap \term{K}^{?_1 \to ?_2 \to ?_1}}_{?_4} \ap \term{I}^{?_3 \to ?_3}}_{?_5} \ap x^{?_0}}_{?_6} \\
            ?_0 = (?_1 \to ?_2 \to ?_1) \to ?_4 \\
            ?_4 = (?_3 \to ?_3) \to ?_5 \\
            ?_5 = ?_0 \to ?_6 \\
            \tau = ?_0 \to ?_6
        \end{gather}
        Подставим все что можно:
        \begin{gather}
            ?_0 = (?_1 \to ?_2 \to ?_1) \to (?_3 \to ?_3) \to ?_0 \to ?_6 \\
            \tau = ?_0 \to ?_6
        \end{gather}
        \(?_0\) встречается как в левой так и в правой части уравнения, значит не существует конечного типа который можно подставить в качестве \(?_0\), значит весь этот терм не типизируется
    \end{proof}

    \paragraph{} $\lambda f~g~x~y\ldotp f \ap g \ap (g \ap x) \ap (g \ap y)$

    \begin{proof}
        \begin{gather}
            \lambda f^{?_0}~g^{?_1}~x^{?_3}~y^{?_6}\ldotp \underbrace{\underbrace{\overbrace{f \ap g}^{?_2} \ap (\overbrace{g \ap x}^{?_4})}_{?_5} \ap (\overbrace{g \ap y}^{?_7})}_{?_8} \\
            \begin{cases}
                ?_0 = ?_1 \to ?_2 \\
                ?_1 = ?_3 \to ?_4 \\
                ?_2 = ?_4 \to ?_5 \\
                ?_1 = ?_6 \to ?_7 \\
                ?_5 = ?_7 \to ?_8 \\
                \tau = ?_0 \to ?_1 \to ?_3 \to ?_6 \to ?_8
            \end{cases}
            \begin{cases}
                ?_0 = (?_3 \to ?_4) \to ?_2 \\
                ?_2 = ?_4 \to ?_5 \\
                ?_6 = ?_3 \\
                ?_7 = ?_4 \\
                ?_5 = ?_7 \to ?_8 \\
                \tau = ?_0 \to (?_3 \to ?_4) \to ?_3 \to ?_6 \to ?_8
            \end{cases}
            \begin{cases}
                ?_0 = (?_3 \to ?_4) \to ?_4 \to ?_5 \\
                ?_5 = ?_4 \to ?_8 \\
                \tau = ?_0 \to (?_3 \to ?_4) \to ?_3 \to ?_3 \to ?_8
            \end{cases} \\
            \begin{cases}
                \tau = ((?_3 \to ?_4) \to ?_4 \to ?_4 \to ?_8) \to (?_3 \to ?_4) \to ?_3 \to ?_3 \to ?_8
            \end{cases}
        \end{gather}
        $\lambda f~g~x~y\ldotp f \ap g \ap (g \ap x) \ap (g \ap y) : ((\alpha \to \beta) \to \beta \to \beta \to \gamma) \to (\alpha \to \beta) \to \alpha \to \alpha \to \gamma$
    \end{proof}

    \paragraph{} $\lambda x~y~z\ldotp x~(y~z~(z~x))$

    \begin{proof}
        \begin{gather}
            \lambda x^{?_0}~y^{?_3}~z^{?_4}\ldotp \overbrace{x~(\underbrace{\overbrace{y~z}^{?_5}~(\overbrace{z~x}^{?_6})}_{?_1})}^{?_2} \\
            \begin{cases}
                ?_0 = ?_1 \to ?_2 \\
                ?_3 = ?_4 \to ?_5 \\
                ?_4 = ?_0 \to ?_6 \\
                ?_5 = ?_6 \to ?_1 \\
                \tau = ?_0 \to ?_3 \to ?_4 \to ?_2
            \end{cases}
            \begin{cases}
                ?_0 = ?_1 \to ?_2 \\
                ?_3 = (?_0 \to ?_6) \to ?_6 \to ?_1 \\
                \tau = ?_0 \to ?_3 \to (?_0 \to ?_6) \to ?_2
            \end{cases}
            \begin{cases}
                ?_3 = ((?_1 \to ?_2) \to ?_6) \to ?_6 \to ?_1 \\
                \tau = (?_1 \to ?_2) \to ?_3 \to ((?_1 \to ?_2) \to ?_6) \to ?_2
            \end{cases} \\
            \begin{cases}
                \tau = (?_1 \to ?_2) \to (((?_1 \to ?_2) \to ?_6) \to ?_6 \to ?_1) \to ((?_1 \to ?_2) \to ?_6) \to ?_2
            \end{cases}
        \end{gather}
        $\lambda x~y~z\ldotp x~(y~z~(z~x)) : (\alpha \to \beta) \to (((\alpha \to \beta) \to \gamma) \to (\gamma \to \alpha)) \to ((\alpha \to \beta) \to \gamma) \to \beta$
    \end{proof}

    \subsubsection{(1б)}

    Возьмите ваш любимый типизированный язык программирования (да, в питоне тоже можно писать типы) и запишите на нём комбинатор $\term{S}$.
    Сообщите ему наиболее общий тип из возможных.
    \begin{itemize}
        \item Проверьте в тесте, что ваш комбинатор справляется с наиболее общим возможным случаем типов его аргументов.
        \item Проверьте, что ваш код проходит проверку типов (если питон, то с помощью mypy).
        \item Указание: для реализации воспользуйтесь в вашем языке концепцией дженериков (generics).
        В C++ она выражается в виде шаблонов (templates).
    \end{itemize}

    \begin{proof}
        \[ S \equiv \lambda x~y~z\ldotp x~z~(y~z) : (\alpha \to \beta \to \gamma) \to (\alpha \to \beta) \to \alpha \to \gamma \]
\begin{minted}[frame=lines,linenos=true,mathescape]{swift}
func S<A, B, C>(x: (A) -> (B) -> C): ((A) -> B) -> (A) -> C {
    { y: (A) -> B => { z: A => x(z)(y(z)) } }
}

main() {
    let unit = S<String, Int, Unit> { _: String => { _: Int => () } }({ _: String => 0 })("")
    println(unit)
}
\end{minted}
    \end{proof}

    \subsubsection{}

    Типизируйте по Чёрчу заданные термы.
    Под типизацией по Чёрчу подразумевается
    расстановка явным образом типовых аннотаций, то есть расписывание термов в стиле
    Чёрча.
    Например, терм $\term{I} \ap \term{K} \ap \term{I}$ типизируется по Чёрчу
    так:

    \[((\lambda x^{(\alpha \to \alpha) \to \beta \to (\alpha \to \alpha)}\ldotp x) \ap (\lambda x^{\alpha \to \alpha} \ap y^\beta\ldotp x) \ap (\lambda x^\alpha\ldotp x)):
    \beta \to \alpha \to \alpha\]

    \paragraph{(1б)} $\term{S} \ap \term{K} \ap \term{I}$

    \begin{proof}
        \begin{gather}
            (\lambda x^{?_0}~y^{?_1}~z^{?_2}\ldotp x~z~(\underbrace{y~z}_{?_6})) \ap (\lambda x^{?_3}~y^{?_4}\ldotp x) \ap (\lambda x^{?_5}\ldotp x) \\
            \begin{cases}
                ?_0 = ?_3 \to ?_4 \to ?_3 \\
                ?_1 = ?_5 \to ?_5 \\
                ?_1 = ?_2 \to ?_6 \\
                ?_0 = ?_2 \to ?_6 \to ?_7
            \end{cases}
            \begin{cases}
                ?_3 = ?_2 \\
                ?_4 = ?_6 \\
                ?_3 = ?_7 \\
                ?_5 = ?_2 \\
                ?_5 = ?_6 \\
                ?_0 = ?_2 \to ?_6 \to ?_7 \\
                ?_1 = ?_2 \to ?_6
            \end{cases}
            \begin{cases}
                ?_0 = ?_2 \to ?_2 \to ?_2 \\
                ?_1 = ?_2 \to ?_2 \\
                ?_3 = ?_2 \\
                ?_4 = ?_2 \\
                ?_5 = ?_2
            \end{cases}
        \end{gather}
        \( (\lambda x^{\alpha \to \alpha \to \alpha}~y^{\alpha \to \alpha}~z^\alpha\ldotp x~z~(y~z)) \ap (\lambda x^\alpha~y^\alpha\ldotp x) \ap (\lambda x^\alpha\ldotp x) : \alpha \to \alpha \)
    \end{proof}

    \paragraph{(1б)} $\term{S} \ap \term{K} \ap \term{K}$

    \begin{proof}
        \begin{gather}
            (\lambda x^{?_0}~y^{?_1}~z^{?_2}\ldotp \overbrace{\underbrace{x~z}_{?_8}~(\underbrace{y~z}_{?_7})}^{?_9}) \ap (\lambda x^{?_3}~y^{?_4}\ldotp x) \ap (\lambda x^{?_5}~y^{?_6}\ldotp x) \\
            \begin{cases}
                ?_0 = ?_3 \to ?_4 \to ?_3 \\
                ?_0 = ?_2 \to ?_8 \\
                ?_1 = ?_5 \to ?_6 \to ?_5 \\
                ?_1 = ?_2 \to ?_7 \\
                ?_8 = ?_7 \to ?_9
            \end{cases}
            \begin{cases}
                ?_0 = ?_2 \to ?_4 \to ?_2 \\
                ?_1 = ?_2 \to ?_6 \to ?_2 \\
                ?_2 = ?_3 \\
                ?_5 = ?_2 \\
                ?_7 = ?_6 \to ?_2 \\
                ?_8 = ?_4 \to ?_2 \\
                ?_8 = ?_7 \to ?_9
            \end{cases}
            \begin{cases}
                ?_0 = ?_2 \to (?_6 \to ?_2) \to ?_2 \\
                ?_1 = ?_2 \to ?_6 \to ?_2 \\
                ?_2 = ?_3 \\
                ?_5 = ?_2 \\
                ?_4 = ?_6 \to ?_2
            \end{cases} \\
            \begin{cases}
                ?_0 = \alpha \to (\beta \to \alpha) \to \alpha \\
                ?_1 = \alpha \to \beta \to \alpha \\
                ?_2 = \alpha \\
                ?_3 = \alpha \\
                ?_4 = \beta \to \alpha \\
                ?_5 = \alpha \\
                ?_6 = \beta
            \end{cases}
        \end{gather}
        \( (\lambda x^{\alpha \to (\beta \to \alpha) \to \alpha}~y^{\alpha \to \beta \to \alpha}~z^\alpha\ldotp x~z~(y~z)) \ap (\lambda x^\alpha~y^{\beta \to \alpha}\ldotp x) \ap (\lambda x^\alpha~y^\beta\ldotp x) : \alpha \to \alpha \)
    \end{proof}

    \subsection{Население типов}

    Локальный минимум: 3б.

    Населите как можно большим количеством замкнутых термов типы или покажите, как
    сконструировать бесконечное их число ($\alpha\beta\eta$-эквивалентные
    термы перечислять не нужно).

    \subsubsection{(1б)}

    \paragraph{} $\alpha \to \alpha$

    \begin{proof}
        $\lambda x\ldotp x$
    \end{proof}

    \paragraph{} $\alpha \to \alpha \to \alpha$

    \begin{proof}
        \begin{itemize}
            \item $\lambda x~y\ldotp x$
            \item $\lambda x~y\ldotp y$
        \end{itemize}
    \end{proof}

    \paragraph{} $\alpha \to \alpha \to \alpha \to \alpha$

    \begin{proof}
        \begin{itemize}
            \item $\lambda x~y~z\ldotp x$
            \item $\lambda x~y~z\ldotp y$
            \item $\lambda x~y~z\ldotp z$
        \end{itemize}
    \end{proof}

    \paragraph{} $\alpha$

    \begin{proof}
        Не существует, т.к. нельзя использовать абстракцию потому что тогда будет стрелочный тип и контекст пустой, т.е. не откуда взять никакой терм типа \(\alpha\)
    \end{proof}

    \subsubsection{(1б)}

    \paragraph{} $\alpha \to \beta \to \alpha$

    \begin{proof}
        $\lambda x~y\ldotp x$
        %% FIXME(iliayar): More?
    \end{proof}

    \paragraph{} $\alpha \to \beta \to (\alpha \to \beta \to \gamma) \to \gamma$

    \begin{proof}
        $\lambda x~y~f\ldotp f \ap x \ap y$
    \end{proof}

    \subsubsection{(1б)}

    \paragraph{} $(\alpha \to \alpha) \to \alpha \to \alpha$

    \begin{proof}
        \[ \lambda f~x\ldotp x \]
        Это числа Черча. Можем бесконечно много раз заменять терм $x$ на терм $f \ap x$:
        \begin{align}
            & \lambda f~x\ldotp x \\
            & \lambda f~x\ldotp f \ap x \\
            & \lambda f~x\ldotp f \ap (f \ap x)
        \end{align}
    \end{proof}

    \paragraph{} $(\alpha \to \beta \to \beta) \to \beta \to \beta$

    \begin{proof}
        $\lambda f~x\ldotp x$
    \end{proof}

    \subsubsection{}

    За одного обитателя даётся половина баллов задания, за бесконечно много обитателей --- полный балл.

    \paragraph{(1б)} $((\alpha \to \beta) \to \alpha) \to (\alpha \to \alpha \to \beta) \to \beta$

    \begin{proof}
        \[ \lambda f^{(\alpha \to \beta) \to \alpha}~g^{\alpha \to \alpha \to \beta}\ldotp g \ap (f \ap (\lambda x\ldotp g \ap x \ap x)) \ap (f \ap (\lambda x\ldotp g \ap x \ap x)) \]
        Можно получить "новое"\ значение типа \(\alpha\) так: \(f \ap (\lambda x\ldotp g \ap x \ap x)\). Можем подставлять вместо перменной \(x\) внутри этого терма его же и получать новый терм. Этот терм не редуцируется, поэтому весь терм будет отличным от предыдущего.
        \begin{align}
            \lambda f~g\ldotp g \ap (f \ap (\lambda x\ldotp g & \ap x \ap x)) \ap (f \ap (\lambda x\ldotp g \ap x \ap x)) \\
            \lambda f~g\ldotp g \ap (f \ap (\lambda x\ldotp g & \ap (f \ap (\lambda x\ldotp g \ap x \ap x)) \ap x)) \ap (f \ap (\lambda x\ldotp g \ap x \ap x)) \\
            \lambda f~g\ldotp g \ap (f \ap (\lambda x\ldotp g & \ap (f \ap (\lambda x\ldotp g \ap x \ap x)) \ap (f \ap (\lambda x\ldotp g \ap x \ap x)))) \ap (f \ap (\lambda x\ldotp g \ap x \ap x))
        \end{align}
    \end{proof}

    \paragraph{(1б)} $((((\alpha \to \beta) \to \alpha) \to \alpha) \to \beta) \to \beta$

    \begin{proof}
        \[ \lambda f^{(((\alpha \to \beta) \to \alpha) \to \alpha) \to \beta}\ldotp f \ap \lambda g^{(\alpha \to \beta) \to \alpha}\ldotp g \ap \lambda x^\alpha\ldotp f \ap \lambda g'\ldotp x \]
        Можем получить "новое"\ значение \(x'\), из \(g'\), т.е. подставив вместо \(x\) терм \(g' \ap \lambda x'\ldotp f \ap \lambda g''\ldotp f \ap x'\). Можем получить бесконечное количество термов делая замену \(x^{(n)}\) на \(g^{(n + 1)} \ap \lambda x^{(n + 1)}\ldotp f \ap \lambda g^{(n + 1)}\ldotp f \ap x^{(n + 1)}\).

        \begin{align}
            & \lambda f\ldotp f \ap \lambda g\ldotp g \ap \lambda x\ldotp f \ap \lambda g'\ldotp x \\
            & \lambda f\ldotp f \ap \lambda g\ldotp g \ap \lambda x\ldotp f \ap \lambda g'\ldotp g' \ap \lamda x'\ldotp f \ap \lambda g''\ldotp f \ap x' \\
            & \lambda f\ldotp f \ap \lambda g\ldotp g \ap \lambda x\ldotp f \ap \lambda g'\ldotp g' \ap \lamda x'\ldotp f \ap \lambda g''\ldotp f \ap g'' \ap \lamda x''\ldotp f \ap \lambda g'''\ldotp f \ap x''
        \end{align}
    \end{proof}

    \subsection{Построение деревьев вывода}

    Локальный минимум: 2б.

    Пример дерева с использованием пакета \texttt{proof}:
    %! suppress = EscapeAmpersand
    \[\resizebox{0.8\hsize}{!}{
        \infer[intro\to]
        {\emptyset \vdash \lambda f~g~x.~f~(g~x) : (\beta \to \gamma) \to (\alpha \to \beta) \to \alpha \to \gamma}
        {
            \infer[intro\to]
            {\{f: \beta \to \gamma\} \vdash \lambda g~x.~f~(g~x) : (\alpha \to \beta) \to \alpha \to \gamma}
            {
                \infer[intro\to]
                {\{f : \beta \to \gamma, g : \alpha \to \beta\} \vdash \lambda x.~f~(g~x) : \alpha \to \gamma}
                {
                    \infer[elim\to]{\{f : \beta \to \gamma, g : \alpha \to \beta, x:  \alpha\} \vdash f~(g~x) : \gamma}
                    {
                        \infer[ctx]{\Gamma \vdash f : \beta \to \gamma}{} &
                        \infer[elim\to]{\Gamma \vdash g~x : \beta}{
                            \infer[ctx]{\Gamma \vdash g : \alpha \to \beta}{} &
                            \infer[ctx]{\Gamma \vdash x : \alpha}{}
                        }
                    }
                }
            }
        }
    }\]

    \subsubsection{(2б)}

    Постройте дерево вывода типа $(\alpha \to \beta \to \gamma) \to (\alpha \to \beta) \to \alpha \to \gamma$ для комбинатора $\term{S}$ в стиле Карри.
    В ответе достаточно одного дерева, промежуточные шаги не нужны.

    \begin{proof}
      \[\resizebox{0.8\hsize}{!}{
          \infer[intro\to]
          {\emptyset \vdash \lambda f~g~x\ldotp f \ap x \ap (g \ap x) : (\alpha \to \beta \to \gamma) \to (\alpha \to \beta) \to \alpha \to \gamma}
          {
              \infer[intro\to]{
                  \{f : \alpha \to \beta \to \gamma\} \vdash \lambda g~x\ldotp f \ap x \ap (g \ap x) : (\alpha \to \beta) \to \alpha \to \gamma        
              }{
                \infer[intro\to]{
                    \{f : \alpha \to \beta \to \gamma, g : \alpha \to \beta\} \vdash \lambda x\ldotp f \ap x \ap (g \ap x) : \alpha \to \gamma        
                }{
                  \infer[elim\to]{
                      \{f : \alpha \to \beta \to \gamma, g : \alpha \to \beta , x : \alpha\} \vdash f \ap x \ap (g \ap x) : \gamma
                  }{
                    \infer[elim\to]{
                        \Gamma \vdash f \ap x : \beta \to \gamma
                    }{
                        \infer[ctx]{
                            \Gamma \vdash f : \alpha \to \beta \to \gamma
                        }{} &
                        \infer[ctx]{
                            \Gamma \vdash x : \alpha
                        }{}
                    } &
                    \infer[elim\to]{
                        \Gamma \vdash g \ap x : \beta
                    }{
                        \infer[ctx]{
                            \Gamma \vdash g : \alpha \to \beta
                        }{} &
                        \infer[ctx]{
                            \Gamma \vdash x : \alpha
                        }{}
                    }
                  }
                }
              }
          }
      }\]
    \end{proof}

\end{document}
